\paragraph{Residual boards and their predictor~\cite[\S2 and \S6.1]{tc}.} The \emph{ordinary board} is the
usual \(n\times n\) board, empty at the start. A \emph{residual board} \(B\) is what a partial placement of
nonattacking queens leaves of it: it keeps the original coordinates and consists of the \(m\) empty rows
and the \(m\) empty columns. A diagonal of either family is \emph{free} if it contains no queen. The
\emph{support} \(E(B)\) is the set of squares that lie in an empty row and an empty column and whose two
diagonals are free; these are the available squares. We write \(Q(B)\) for the number of solutions that
extend the placement and \(B+e\) for the child obtained by placing a queen at \(e\in E(B)\). The initial
board \(B_n\) is the residual board of the empty placement, so \(Q(B_n)=Q(n)\). For weights
\(p=(p_e)_{e\in E(B)}\) and a free diagonal \(f\), let \(\ell_f=\sum_{e\in f}p_e\) and \(h_f=1-\ell_f\).
The Bethe entropy of \(B\) is
\begin{equation}\label{eq:bethe-residual}
 H_B(p)=-\sum_e p_e\log p_e+3\sum_e(1-p_e)\log(1-p_e)-\sum_f h_f\log h_f
\end{equation}
on the domain \(0<p_e<1/4\), \(h_f>0\), with unit row and column masses. There \(-H_B\) is strictly
convex, so \(H_B\) has at most one stationary point, the \emph{Bethe point} of \(B\). Let \(\mathcal A\) be
the incidence matrix of the squares of \(E(B)\) with the lines of \(B\), namely its empty rows, its empty
columns, and the free diagonals that meet \(E(B)\), and let \(\tau_4(u)=u(1-u)/(1-4u)\). The line
operator is
\[
 S_B=\operatorname{diag}(0_R,0_C,h)+\mathcal A^{\mathsf T}\operatorname{diag}(\tau_4(p_e))\mathcal A .
\]
When the bipartite graph of rows and columns of the support is connected, its kernel is spanned by
\((\mathbf1_R,-\mathbf1_C,0)\). The predictor of \(B\) is then
\begin{equation}\label{eq:predictor-residual}
 L(B)=H_B(p)-\tfrac12\sum_e\bigl[\log(1-4p_e)-4\log(1-p_e)\bigr]-\tfrac12\log\operatorname{pdet}S_B
      +\tfrac12\log(4\pi m),
\end{equation}
evaluated at the Bethe point, where \(\operatorname{pdet}\) is the product of the positive eigenvalues.
Deleting from \(S_B\) the row and column of one axial line gives a cofactor equal to
\(\operatorname{pdet}S_B/(2m)\); used with \(\tfrac12\log(2\pi)\), as in~\cite[\S\S6.1 and 7.1]{pc}, it
gives the same \(L(B)\). On the initial board and for large \(n\), \(L(B_n)=\log Q_{\rm pred}(n)\), the
predictor of Appendix~\ref{app:predictor}.

\paragraph{Histories and the two laws~\cite[\S2]{tc}.} Draw a uniformly random order \(\sigma\) of the
rows. Under the uniform law \(\mu\), a uniform solution is chosen independently of \(\sigma\) and its
queens are revealed in that order. Under the adaptive law \(\nu\), given the next row \(i\), with squares
\(R_i\), the square \(e\in R_i\cap E(B)\) is chosen with probability
\[
 \nu_B(e\mid i)=\frac{\exp L(B+e)}{\sum_{e'\in R_i\cap E(B)}\exp L(B+e')},\qquad
 \rho(B,i)=\sum_{e'\in R_i\cap E(B)}\exp\{L(B+e')-L(B)\}.
\]
A \emph{history} is the resulting sequence of boards, together with the row order. A board with \(m\)
empty rows has size \(m\).

\paragraph{The stopping rule~\cite[Def.~2.2]{tc}.} Fix \(0<\delta<1\), \(M=\lceil\delta n\rceil\), and an
integer \(K\) with \(K_0\le K\le M\). The threshold \(K_0\), the constants, and the tolerances are those
of~\cite[\S5.4]{tc}. Let \(B\) be a board of a history, of size \(m\).
\begin{enumerate}
\item[(a)] \emph{Macroscopic range}, \(m\ge M\). The board is \emph{regular} if the discrepancy of the
history on reaching it, defined in~\cite[\S4.2]{tc} with the partition of step \(a_{\rm mac}\)
of~\cite[Prop.~4.11]{tc}, is smaller than the radius \(\varrho_{\rm mac}\) of that proposition.
\item[(b)] \emph{Cut}, \(m=M\). In addition, the observables of~\cite[\S5.3]{tc} lie within their entry
tolerances and \(D_{\max}(M)<\alpha_{\rm hist}M\).
\item[(c)] \emph{Mesoscopic range}, \(K\le m<M\). The board is regular if its Bethe point exists and
belongs to the inner class: it satisfies the diffusion, vacancy, and gap
bounds~\eqref{eq:clase-importada} with the constants fixed in~\cite[\S5.4]{tc}, and the mass of each
diagonal is at most \(\epsilon\).
\end{enumerate}
The \emph{first exit} of a history is its first board that is not regular; its parent is regular. The
history stops at the first exit or, if there is none, on reaching size \(K\). The persistence event is
that there is no exit: all boards down to size \(K\), inclusive, are regular. In this paper the floor is
\(K=j=\lfloor bn\rfloor\) (Section~\ref{sec:signedmap}).
